% TOP_PROBLEM: {"id": 5500037, "title": "Counting Polyominoes", "category_id": 2, "category": {"id": 2, "name": "combinatorics", "display_name": "Combinatorics", "description": "Counting problems, graph theory, discrete structures.", "slug": "combinatorics", "order_index": 2, "created_at": "2026-07-31T15:26:25.671Z"}, "status": "open", "problem_number": "AMR-054-0037", "set_id": 13}
% FIRST_POSTED: 2026-10-01
% ENTRY_KIND: statement_only
% UnsolvedMath ID 5500037; source catalogue code AMR-054-0037
% !TeX program = lualatex
% Definition and sources only; this file is not an LLM solution attempt.
\documentclass[11pt,a4paper]{article}
\usepackage[margin=25mm]{geometry}
\usepackage{fontspec}
\setmainfont{lmroman10-regular.otf}[BoldFont=lmroman10-bold.otf,
 ItalicFont=lmroman10-italic.otf,BoldItalicFont=lmroman10-bolditalic.otf]
\usepackage{amsmath,amssymb,mathtools}
\usepackage{unicode-math}
\setmathfont{latinmodern-math.otf}
\AtBeginDocument{\let\setminus\smallsetminus}
\usepackage{xurl,hyperref}
\hypersetup{colorlinks=true,urlcolor=blue}
\setcounter{secnumdepth}{0}
\setlength{\parindent}{0pt}
\setlength{\parskip}{5pt}
\setlength{\emergencystretch}{3em}
\begin{document}
\section*{Counting Polyominoes}\label{5500037}
{\small UnsolvedMath ID 5500037 · AMR-054-0037 · Combinatorics · AMR Open Problem Lists}
\subsection{Definitions and mathematical statement}
A polyomino of order $n$ is the union of $n$ distinct unit squares
$[i,i+1]\times[j,j+1]$, with $(i,j)\in\mathbb Z^2$, such that their
edge-adjacency graph is connected. Meeting at a corner alone does not
give adjacency. Holes are permitted unless an additional restriction
is explicitly imposed.

Let $t_n$ count these objects up to translation only. Let $c_n$ count
them up to translations and rotations through multiples of $\pi/2$,
and let $f_n$ also identify reflections. These are respectively the
fixed, one-sided (called chiral in the source), and free conventions.
Determine exact general counting formulas and sharp asymptotics for
these sequences.

In particular, determine the growth constant
\[
                  \lambda=\lim_{n\to\infty} t_n^{1/n},
\]
and the finer subexponential behavior of $t_n$. The limit is a growth
rate per added square: its defining relation is
$\log t_n/n\to\log\lambda$, not a polynomial-growth claim
$t_n\approx n^\lambda$.
Enumerating finitely many values does not determine this asymptotic
quantity. The conventions also matter for exact counts: symmetric
polyominoes have smaller rotation or reflection orbits, so the three
sequences cannot be obtained by dividing every term by four or eight.
\subsection{Short English statement}
Count connected collections of grid squares, distinguishing rotations
and reflections correctly, and determine their exponential growth rate
and finer asymptotics.
\subsection{Sources}
\begin{itemize}
\item[S1] ULAM.AI and the UnsolvedMath Contributors, supplied problems.json, record 5500037 (AMR-054-0037), AMR Open Problem Lists. Catalogue identity and source transcription; CC BY 4.0.
\url{https://huggingface.co/datasets/ulamai/UnsolvedMath}.
\item[S2] Open Problems Project - Computational Geometry, Problem 37. Original source indexed by AMR-054-0037.
\url{https://topp.openproblem.net/p37}.
\end{itemize}
\end{document}
